\input{preamble}
\usepackage{amsthm}
\usepackage{algorithm}
\usepackage{algpseudocode}
\newtheorem{lemma}{Lemma}
\newtheorem{theorem}{Theorem}
\newtheorem{remark}{Remark}

\title{Problem 10 (Lemmas L1, L3, L5): SPD operator, structured mat--vec, and exact complexity}
\date{}

\begin{document}
\maketitle

\noindent
This note proves the spectral properties of the mode-$k$ normal-equation
operator that are the prerequisite for using the conjugate gradient (CG)
method (Lemma~\ref{lem:psd}--\ref{lem:singular}), the structured
matrix--vector product used inside CG (Lemma~\ref{lem:matvec}), and the exact
complexity of the preconditioned CG (PCG) solver
(Lemma~\ref{lem:precond}--\ref{lem:complexity}).  Throughout, with the notation
of the problem statement, set
\begin{equation}\label{eq:Adef}
  \aA \;=\; (Z\otimes K)^{T} S\, S^{T} (Z\otimes K)\;+\;\lambda\,(I_r\otimes K)
  \;\in\;\bR^{nr\times nr},
\end{equation}
and let
\begin{equation}\label{eq:rhs}
  b \;=\; (I_r\otimes K)\,\vecop(B),\qquad B = TZ,
\end{equation}
so that the subproblem is $\aA\,\vecop(W)=b$. We always assume, as in the
problem statement, that $K=K^{T}\in\bR^{n\times n}$ is positive semidefinite
($K\succeq0$, the RKHS kernel matrix), that $S\in\bR^{N\times q}$ is a column
subset of $I_N$ (hence $S^{T}S=I_q$ and $SS^{T}$ is an orthogonal projector),
and that $\lambda\ge0$.  We use $\odot$ for the column-wise Khatri--Rao product
and $\circledast$ for the Hadamard (entrywise) product.

\bigskip

\begin{lemma}[Symmetry and positive semidefiniteness]\label{lem:psd}
For every $K\succeq0$ and every $\lambda\ge0$ the operator $\aA$ in
\eqref{eq:Adef} is symmetric and positive semidefinite, i.e.
$\aA=\aA^{T}$ and $x^{T}\aA\,x\ge0$ for all $x\in\bR^{nr}$.
\end{lemma}

\begin{proof}
\emph{Symmetry.} Write $G:=(Z\otimes K)^{T} S\,S^{T} (Z\otimes K)$ and
$R:=\lambda(I_r\otimes K)$, so $\aA=G+R$.
For any matrix $M:=S^{T}(Z\otimes K)\in\bR^{q\times nr}$ we have
$G=M^{T}M$, which is symmetric because $(M^{T}M)^{T}=M^{T}M$.
For $R$, the transpose of a Kronecker product satisfies
$(I_r\otimes K)^{T}=I_r^{T}\otimes K^{T}=I_r\otimes K$ because $I_r$ and $K$
are symmetric; multiplying by the scalar $\lambda$ preserves symmetry, so
$R^{T}=R$. Hence $\aA^{T}=G^{T}+R^{T}=G+R=\aA$.

\smallskip
\emph{Positive semidefiniteness.} Fix $x\in\bR^{nr}$. Then
\begin{equation}\label{eq:quadform}
  x^{T}\aA\,x
  = x^{T}M^{T}M x + \lambda\, x^{T}(I_r\otimes K)x
  = \underbrace{\|Mx\|_2^{2}}_{\ge0}
   \;+\; \lambda\, x^{T}(I_r\otimes K)x .
\end{equation}
The first term is a squared Euclidean norm, hence $\ge0$. For the second term,
since $K\succeq0$ there is a factorization $K=L L^{T}$ (e.g.\ a Cholesky or
symmetric square-root factor $L\in\bR^{n\times n}$). Using the Kronecker
identity $I_r\otimes K=I_r\otimes(LL^{T})=(I_r\otimes L)(I_r\otimes L)^{T}$,
\[
  x^{T}(I_r\otimes K)x
  = x^{T}(I_r\otimes L)(I_r\otimes L)^{T}x
  = \bigl\|(I_r\otimes L)^{T}x\bigr\|_2^{2}\;\ge\;0 .
\]
Because $\lambda\ge0$, the second term in \eqref{eq:quadform} is $\ge0$, and
therefore $x^{T}\aA x\ge0$. Thus $\aA\succeq0$.
\end{proof}

\begin{lemma}[Strict positive definiteness]\label{lem:spd}
Assume in addition that $K\succ0$ (i.e.\ $K$ is nonsingular) and $\lambda>0$.
Then $\aA$ is symmetric \emph{positive definite}: $x^{T}\aA x>0$ for all
$x\neq0$, equivalently $\aA$ is nonsingular. These two conditions are also
\emph{necessary} in general: if $\lambda=0$ then $\aA$ is singular, and if
$K$ is singular then $\aA$ is singular.
\end{lemma}

\begin{proof}
\emph{Sufficiency.} Symmetry holds by Lemma~\ref{lem:psd}. Let $x\neq0$.
By \eqref{eq:quadform},
$x^{T}\aA x \ge \lambda\,x^{T}(I_r\otimes K)x$. Since $K\succ0$, its
eigenvalues are bounded below by $\mu_{\min}(K)>0$, and the eigenvalues of the
Kronecker product $I_r\otimes K$ are exactly the eigenvalues of $K$, each with
multiplicity $r$ (a standard fact: if $K v=\mu v$ then
$(I_r\otimes K)(e_j\otimes v)=\mu\,(e_j\otimes v)$ for each standard basis
vector $e_j$, giving $nr$ orthonormal eigenpairs in total). Hence
$I_r\otimes K\succeq \mu_{\min}(K)\,I_{nr}$ and
\[
  x^{T}\aA x \;\ge\; \lambda\,\mu_{\min}(K)\,\|x\|_2^{2} \;>\;0 ,
\]
because $\lambda>0$, $\mu_{\min}(K)>0$, and $x\neq0$. A symmetric operator
with strictly positive quadratic form on all nonzero vectors is positive
definite and in particular nonsingular.

\smallskip
\emph{Necessity of $\lambda>0$.} If $\lambda=0$, then
$\aA = (Z\otimes K)^{T}SS^{T}(Z\otimes K) = M^{T}M$ with
$M=S^{T}(Z\otimes K)\in\bR^{q\times nr}$. Since $M$ has only $q$ rows and
$\aA$ has order $nr$, whenever $q<nr$ the rank of $\aA$ is at most $q<nr$, so
$\aA$ is singular.

\smallskip
\emph{Necessity of $K\succ0$.} Suppose $K$ is singular, so there is
$v\neq0$ with $Kv=0$. Pick any $e_j$ ($j\in\{1,\dots,r\}$) and set
$x:=e_j\otimes v\neq0$. Using the mixed-product rule
$(A\otimes B)(C\otimes D)=(AC)\otimes(BD)$,
\[
  (Z\otimes K)x=(Ze_j)\otimes(Kv)=0,
  \qquad
  (I_r\otimes K)x=(I_r e_j)\otimes(Kv)=0 .
\]
Hence $\aA x=0$ for any $\lambda$, so $x\in\ker\aA$ and $\aA$ is singular.
\end{proof}

\begin{lemma}[Consistency in the singular--$K$ boundary case]
\label{lem:singular}
Let $K\succeq0$ be (possibly) singular and $\lambda>0$. Then
$\ker\aA=\ker(I_r\otimes K)$, hence
$\operatorname{range}(\aA)=\operatorname{range}(I_r\otimes K)$, and the
right-hand side $b=(I_r\otimes K)\vecop(B)\in\operatorname{range}(\aA)$, so
$\aA w=b$ is consistent; CG started at $w_0=0$ converges to the minimum-norm
solution.
\end{lemma}

\begin{proof}
Write $\aA=M^{T}M+\lambda(I_r\otimes K)$ with $M=S^{T}(Z\otimes K)$; both
summands are PSD. For PSD $P,Q$, $(P+Q)x=0$ gives
$0=x^{T}Px+x^{T}Qx$ with both terms $\ge0$, forcing $Px=Qx=0$, so
$\ker(P+Q)=\ker P\cap\ker Q$. Here
$\ker\aA=\ker M\cap\ker(I_r\otimes K)$. Moreover
$(Z\otimes K)=(Z\otimes I_n)(I_r\otimes K)$, so $(I_r\otimes K)x=0$ implies
$Mx=0$; thus $\ker(I_r\otimes K)\subseteq\ker M$ and
$\ker\aA=\ker(I_r\otimes K)$. Symmetry gives
$\operatorname{range}(\aA)=\ker(\aA)^{\perp}=\ker(I_r\otimes K)^{\perp}
=\operatorname{range}(I_r\otimes K)$. Since $b$ lies in the column space of
$I_r\otimes K$, it lies in $\operatorname{range}(\aA)$; the system is
consistent. With $w_0=0$ the residual $r_0=b\in\operatorname{range}(\aA)$ and
the Krylov subspaces stay in the $\aA$-invariant subspace
$\operatorname{range}(\aA)$, so CG behaves as on the SPD restriction
$\aA|_{\operatorname{range}(\aA)}$ and converges to the minimum-norm solution.
\end{proof}

\begin{remark}[Practical safeguard]\label{rem:reg}
When $K\succeq0$ is numerically singular, replace $K$ by
$K_\varepsilon:=K+\varepsilon I_n\succ0$ with small $\varepsilon>0$; then
Lemma~\ref{lem:spd} gives $\aA_\varepsilon\succ0$ for every $\lambda>0$, and as
$\varepsilon\to0^{+}$ the solution of
$\aA_\varepsilon w=(I_r\otimes K_\varepsilon)\vecop(B)$ converges to the
minimum-norm solution of Lemma~\ref{lem:singular}. The shift changes neither
the structure nor the asymptotic cost of the mat--vec or preconditioner below.
Henceforth we assume the SPD regime $K\succ0$, $\lambda>0$ (using
$K_\varepsilon$ if needed).
\end{remark}

\section*{Structured matrix--vector product (Lemma L3)}

The heart of an iterative solver is that $\aA$ is never formed; only the action
$y\mapsto \aA y$ is needed. We give the exact data structures and their cost.

\paragraph{Data extracted from $S$ once.}
Each of the $q$ observed entries corresponds to a position $p\in\{0,\dots,N-1\}$
in $\vecop(T)$, where $T\in\bR^{n\times M}$ is the mode-$k$ unfolding. Writing
$p=a+n\,c$ with $a\in\{0,\dots,n-1\}$ (the mode-$k$ index) and
$c\in\{0,\dots,M-1\}$ (the linear index over the remaining modes), the $j$-th
observation yields a pair $(a_j,c_j)$. From these we read off, \emph{once per
subproblem and without any order-$M$ or order-$N$ work},
\begin{equation}\label{eq:tilde}
  \tilde K \;=\; \bigl[K_{a_j,\,:}\bigr]_{j=1}^{q}\in\bR^{q\times n},
  \qquad
  \tilde Z \;=\; \bigl[Z_{c_j,\,:}\bigr]_{j=1}^{q}\in\bR^{q\times r},
\end{equation}
i.e.\ $\tilde K$ ($\tilde Z$) is the subset of $q$ rows of $K$ (of $Z$) indexed
by the observations, with possible repeats. Reading the $q$ rows of $K$ is
$O(qn)$; row $c_j$ of the Khatri--Rao product $Z$ is the entrywise product of
the corresponding rows of the $d-1$ fixed factors,
$Z_{c_j,:}=\circledast_{i\neq k}A_i\,[\,c_j^{(i)},:\,]$, computable in
$O((d-1)r)$ per row, hence $O(qdr)$ total. Neither $Z\in\bR^{M\times r}$ nor
$\vecop(T)\in\bR^{N}$ is ever materialized.

\begin{lemma}[Structured action of $\aA$]\label{lem:matvec}
Let $y=\vecop(Y)$ with $Y\in\bR^{n\times r}$. Define the column vector
$c\in\bR^{q}$ by
\begin{equation}\label{eq:fwd}
  c_j=\bigl(\tilde K_{j,:}\,Y\bigr)\,\tilde Z_{j,:}^{T}
   =\sum_{a=1}^{n}\sum_{p=1}^{r}\tilde K_{j,a}\,Y_{a,p}\,\tilde Z_{j,p},
   \qquad j=1,\dots,q,
\end{equation}
and set $D:=\diag(c)\in\bR^{q\times q}$. Then
\begin{equation}\label{eq:Aaction}
  \aA\,y
  \;=\;\vecop\!\Bigl(\,\tilde K^{T} D\,\tilde Z \;+\; \lambda\,K Y\,\Bigr).
\end{equation}
No $N$- or $M$-dimensional object appears; the entire computation uses only
$\tilde K\in\bR^{q\times n}$, $\tilde Z\in\bR^{q\times r}$, $K\in\bR^{n\times n}$
and $Y\in\bR^{n\times r}$, and costs
\begin{equation}\label{eq:matveccost}
  \mathrm{cost}(\aA\text{-multiply})
  \;=\; O\!\bigl(qnr + n^{2}r\bigr).
\end{equation}
\end{lemma}

\begin{proof}
\emph{Identity.} By the Kronecker--vec rule
$(Z\otimes K)\vecop(Y)=\vecop(KYZ^{T})$, so the matrix $KYZ^{T}\in\bR^{n\times M}$
is the un-vectorization of $(Z\otimes K)y$. The selector $S^{T}$ extracts the
$q$ observed entries; the $j$-th of these is the entry of $KYZ^{T}$ at row $a_j$,
column $c_j$, namely
\[
  \bigl(S^{T}(Z\otimes K)y\bigr)_j
  =\bigl(KYZ^{T}\bigr)_{a_j,c_j}
  =K_{a_j,:}\,Y\,(Z_{c_j,:})^{T}
  =\tilde K_{j,:}\,Y\,\tilde Z_{j,:}^{T}=c_j,
\]
which is \eqref{eq:fwd}; here we used \eqref{eq:tilde}. Next $SS^{T}$ scatters
$c$ back into the $(a_j,c_j)$ positions of an otherwise-zero $n\times M$ matrix
$E$, i.e.\ $E=\sum_{j}c_j\,e_{a_j}e_{c_j}^{T}$, and
$(Z\otimes K)^{T}SS^{T}(Z\otimes K)y=\vecop(K^{T}EZ)=\vecop(KEZ)$ (again by the
vec rule and $K=K^{T}$). Because $E$ is supported on the $q$ observed positions,
\[
  KEZ=\sum_{j=1}^{q}c_j\,(Ke_{a_j})(e_{c_j}^{T}Z)
     =\sum_{j=1}^{q}c_j\,K_{:,a_j}\,Z_{c_j,:}
     =\tilde K^{T}D\,\tilde Z,
\]
where the last equality is the matrix form of the rank-$q$ sum with
$D=\diag(c)$ and rows \eqref{eq:tilde} (using $K_{:,a_j}=K_{a_j,:}^{T}=\tilde
K_{j,:}^{T}$ by symmetry). Finally
$(I_r\otimes K)y=\vecop(KY)$, so
$\aA y=\vecop(\tilde K^{T}D\tilde Z+\lambda KY)$, which is \eqref{eq:Aaction}.

\smallskip
\emph{Cost.} Forming $\tilde K Y$ ($q\times n$ times $n\times r$) is $O(qnr)$;
the $q$ inner products in \eqref{eq:fwd} cost $O(qr)$; forming $D\tilde Z$
(scaling $q$ rows) is $O(qr)$; $\tilde K^{T}(D\tilde Z)$
($n\times q$ times $q\times r$) is $O(qnr)$; and $\lambda KY$
($n\times n$ times $n\times r$) is $O(n^{2}r)$. Summing gives
$O(qnr+n^{2}r)$, which is \eqref{eq:matveccost}.
\end{proof}

\begin{remark}[Honest mat--vec cost; no hidden structure]\label{rem:honestcost}
The $K$-multiply $Y\mapsto KY$ in \eqref{eq:Aaction} is the only operation
touching the full $n\times n$ kernel; for a \emph{generic dense} $K$ it costs
$\Theta(n^{2}r)$, and we count it as such in \eqref{eq:matveccost}. We do
\emph{not} assume a structured (e.g.\ low-rank Nystr\"om or fast-multipole)
kernel. If $K$ does admit a rank-$\rho$ factorization $K=FF^{T}$
($F\in\bR^{n\times\rho}$), then $KY=F(F^{T}Y)$ costs $O(n\rho r)$ and one would
replace $n^{2}r$ by $n\rho r$ throughout; this is an optional improvement, not
an assumption of the analysis below. Under the standing regime
$n,r<q$, the term $qnr$ dominates $n^{2}r$ (since $q>n$) and $nr^{2}$ (since
$q>r$), so the per-mat--vec cost is $O(qnr)$; we nevertheless keep the
$n^{2}r$ term explicit so the dependence on the dense $K$-multiply is visible.
\end{remark}

\section*{Preconditioner and right-hand side (Lemma L4)}

\begin{lemma}[Exact Kronecker preconditioner and its cost]\label{lem:precond}
Let $G:=Z^{T}Z\in\bR^{r\times r}$ and define the SPD preconditioner
\begin{equation}\label{eq:Pdef}
  P \;:=\; G\otimes K^{2}\;+\;\lambda\,(I_r\otimes K)
  \;=\;(Z\otimes K)^{T}(Z\otimes K)+\lambda(I_r\otimes K),
\end{equation}
i.e.\ $P$ is exactly $\aA$ in the fully observed limit $SS^{T}=I_N$. Then:
\begin{enumerate}[label=\textup{(\roman*)}]
\item $G$ is computable \emph{without order-$M$ work} via the Khatri--Rao
      Gram identity
      \begin{equation}\label{eq:krgram}
        G=Z^{T}Z=\big(A_{d}\odot\!\cdots\!\odot A_{1}\big)^{T}\!\big(A_{d}\odot\!\cdots\!\odot A_{1}\big)
        =\!\!\!\circledast_{\,i\neq k}\!\!\big(A_i^{T}A_i\big),
      \end{equation}
      at cost $O\!\big(r^{2}\sum_{i\neq k} n_i + d\,r^{2}\big)$.
\item With the eigendecompositions $K=U\Theta U^{T}$ ($\Theta=\diag(\theta)$,
      one-time $O(n^{3})$) and $G=V\Gamma V^{T}$ ($\Gamma=\diag(g)$, one-time
      $O(r^{3})$), one has the simultaneous diagonalization
      \begin{equation}\label{eq:Peig}
        P=(V\otimes U)\,\Lambda\,(V\otimes U)^{T},\qquad
        \Lambda=\diag\big(g_p\theta_a^{2}+\lambda\theta_a\big)_{a=1,p=1}^{n,\ r},
      \end{equation}
      and $\Lambda\succ0$ when $K\succ0$, $\lambda>0$. Consequently, for
      $z=\vecop(Z_0)$ with $Z_0\in\bR^{n\times r}$,
      \begin{equation}\label{eq:Pinv}
        P^{-1}z=\vecop\!\Big(U\big[(U^{T}Z_0V)\oslash \Omega\big]V^{T}\Big),
        \qquad \Omega_{a,p}:=g_p\theta_a^{2}+\lambda\theta_a,
      \end{equation}
      ($\oslash$ entrywise division), at cost $O(n^{2}r+nr^{2})$ per apply.
\end{enumerate}
\end{lemma}

\begin{proof}
(i) The mixed-product rule gives
$(Z\otimes K)^{T}(Z\otimes K)=(Z^{T}Z)\otimes(K^{T}K)=G\otimes K^{2}$ (using
$K=K^{T}$), which yields the second equality in \eqref{eq:Pdef}. For
\eqref{eq:krgram}, column $j$ of $Z=A_d\odot\cdots\odot A_1$ is
$\bigotimes_{i\neq k}(A_i)_{:,j}$, so
$(Z^{T}Z)_{jl}=\prod_{i\neq k}\langle (A_i)_{:,j},(A_i)_{:,l}\rangle
=\prod_{i\neq k}(A_i^{T}A_i)_{jl}$, i.e.\ the Hadamard product of the factor
Gram matrices. Each $A_i^{T}A_i$ costs $O(n_i r^{2})$ and the $d-1$ Hadamard
products of $r\times r$ matrices cost $O(d r^{2})$, giving the stated bound; no
factor of $M=\prod_{i\neq k}n_i$ appears.

(ii) Substituting $K=U\Theta U^{T}$, $G=V\Gamma V^{T}$ into \eqref{eq:Pdef} and
using $(V\Gamma V^{T})\otimes(U\Theta^{2}U^{T})=(V\otimes U)(\Gamma\otimes\Theta^{2})(V\otimes U)^{T}$
and $I_r\otimes(U\Theta U^{T})=(V\otimes U)(I_r\otimes\Theta)(V\otimes U)^{T}$
(the latter because $V$ is orthogonal, $V I_r V^{T}=I_r$) gives \eqref{eq:Peig}
with diagonal entries $g_p\theta_a^{2}+\lambda\theta_a$; these are $>0$ when all
$\theta_a>0$ and $\lambda>0$, so $P\succ0$. Inverting the orthogonal factors and
the diagonal yields
$P^{-1}=(V\otimes U)\Lambda^{-1}(V\otimes U)^{T}$, and the vec-rule
$(V\otimes U)\vecop(C)=\vecop(UCV^{T})$ turns this into \eqref{eq:Pinv}.
The cost of one apply is: $U^{T}Z_0$ is $O(n^{2}r)$; $(\cdot)V$ is $O(nr^{2})$;
entrywise division $O(nr)$; left-multiply by $U$ is $O(n^{2}r)$; right by
$V^{T}$ is $O(nr^{2})$; total $O(n^{2}r+nr^{2})$.
\end{proof}

\begin{remark}[The preconditioner is order-$M$ \emph{and} order-$N$ free; the
naive Gram is \emph{not} used]\label{rem:noM}
It is essential to note that $P$ in \eqref{eq:Pdef} is built \emph{without ever
forming $Z\in\bR^{M\times r}$ and without the naive Gram
$Z^{T}Z$}, whose direct evaluation $(r\times M)\cdot(M\times r)$ would cost
$O(Mr^{2})$ --- an order-$M$ operation forbidden by the problem. Instead the
only Gram-type object needed is the $r\times r$ matrix
$G=Z^{T}Z=\circledast_{i\neq k}(A_i^{T}A_i)$, and Lemma~\ref{lem:precond}(i)
computes it from the \emph{factor} Grams $A_i^{T}A_i\in\bR^{r\times r}$ at total
cost
\[
  O\!\Big(r^{2}\sum_{i\neq k}n_i+d\,r^{2}\Big),
\]
which is governed by the \emph{sum} $\sum_{i\neq k}n_i$, never the
\emph{product} $M=\prod_{i\neq k}n_i$. Under the standing hypotheses
$d<n<q$ and $n_i\le\max_i n_i$ one has $\sum_{i\neq k}n_i\le(d-1)\max_i n_i$,
which is $O(dq)$ if all mode sizes obey $n_i<q$, and in any case is
$O\!\big(\sum_i n_i\big)=O(N^{1/d}\,d)$ in the balanced case --- always strictly
sub-$M$. The remaining ingredients of $P$ are the kernel $K$ itself
($n\times n$, available) and its eigendecomposition $K=U\Theta U^{T}$
($O(n^{3})$, one-time); none of these touches an object of size $M$ or $N$.
Thus $P$ satisfies the order-$M$ and order-$N$ bans exactly, while remaining the
\emph{fully-observed} ($SS^{T}=I_N$) operator and hence a tight surrogate for
$\aA$.
\end{remark}

\begin{remark}[Admissible alternatives, including the observed Gram]
\label{rem:altprecond}
Two further order-$M$/$N$-free preconditioners are available and admissible.
\begin{enumerate}[label=\textup{(\alph*)}]
\item \emph{Observed-Gram (sampled) preconditioner.} Replace $G=Z^{T}Z$ by the
      \emph{observed} Gram $\tilde G:=\tilde Z^{T}\tilde Z\in\bR^{r\times r}$,
      assembled from the $q$ observed rows $\tilde Z$ of \eqref{eq:tilde} at cost
      $O(qr^{2})$ (no order-$M$ work, since $\tilde Z$ has only $q$ rows), and
      a diagonal column-weight $W_K:=\diag(\tilde K^{T}\tilde K)$ surrogate for
      $K^{2}$; the simplest admissible choice that mirrors the data term is
      \[
        \tilde P:=\tilde G\otimes K^{2}+\lambda(I_r\otimes K),
      \]
      with the \emph{identical} simultaneous-diagonalization apply of
      Lemma~\ref{lem:precond}(ii) after eigendecomposing $\tilde G=\tilde V\tilde\Gamma\tilde V^{T}$
      ($O(r^{3})$). This is closer to $\aA$ entrywise when the sampling is
      non-uniform, but $\tilde G\succeq0$ may be singular when $q<r$ (excluded
      here since $r<q$) or rank-deficient under adversarial sampling; the shift
      $\lambda(I_r\otimes K)\succ0$ keeps $\tilde P\succ0$ regardless.
\item \emph{Block-Jacobi / regularizer preconditioner.} The bare regularizer
      $P_0:=\lambda(I_r\otimes K)$ is SPD (Lemma~\ref{lem:spd}), order-$M$-free,
      and its apply is $P_0^{-1}z=\lambda^{-1}\vecop(K^{-1}Z_0)$ via the
      one-time Cholesky $K=LL^{T}$ ($O(n^{3})$) at $O(n^{2}r)$ per solve. It is
      the cheapest admissible choice and already bounds the small end of the
      spectrum of $\aA$ (Remark~\ref{rem:whyprecond}); $P$ of \eqref{eq:Pdef}
      additionally captures the data curvature and is preferred.
\end{enumerate}
Each option is SPD for $K\succ0$, $\lambda>0$ by the Kronecker/PSD argument of
Lemma~\ref{lem:precond}(ii), is assembled and applied with no order-$M$ or
order-$N$ operation, and uses only one-time factorizations of size $n$ and $r$.
We adopt $P$ of \eqref{eq:Pdef} as the default because it is the exact
fully-observed operator and gives the spectral sandwich of
Remark~\ref{rem:whyprecond}.
\end{remark}

\begin{lemma}[Spectral equivalence $\lambda(I_r\otimes K)\preceq\aA\preceq P$
and clustering]\label{lem:specequiv}
Let $K\succ0$, $\lambda>0$, and let $P$ be as in \eqref{eq:Pdef}. Then
\begin{equation}\label{eq:sandwich}
  \lambda\,(I_r\otimes K)\;\preceq\;\aA\;\preceq\;P,
\end{equation}
and consequently every generalized eigenvalue $\mu$ of the pencil
$(\aA,P)$ --- equivalently every eigenvalue of the $P$-self-adjoint operator
$P^{-1}\aA$, i.e.\ of the symmetric matrix $P^{-1/2}\aA P^{-1/2}$ --- satisfies
\begin{equation}\label{eq:specbound}
  0\;<\;\kappa_0\;\le\;\mu\;\le\;1,
  \qquad
  \kappa_0:=\min_{x\neq0}\frac{\lambda\,x^{T}(I_r\otimes K)x}{x^{T}Px}\;>\;0 .
\end{equation}
Hence the $P$-preconditioned condition number obeys
$\kappa_P(\aA)=\mu_{\max}/\mu_{\min}\le 1/\kappa_0$, independent of $q$.
\end{lemma}

\begin{proof}
Write $\aA=(Z\otimes K)^{T}SS^{T}(Z\otimes K)+\lambda(I_r\otimes K)$ and
$P=(Z\otimes K)^{T}(Z\otimes K)+\lambda(I_r\otimes K)$. Since $S$ is a column
subset of $I_N$, $SS^{T}$ is an orthogonal projector, so
$0\preceq SS^{T}\preceq I_N$. Conjugating by $(Z\otimes K)$ preserves the
Loewner order: for the lower bound, $0\preceq SS^{T}$ gives
$0\preceq(Z\otimes K)^{T}SS^{T}(Z\otimes K)$, whence
$\lambda(I_r\otimes K)\preceq\aA$; for the upper bound, $SS^{T}\preceq I_N$
gives $(Z\otimes K)^{T}SS^{T}(Z\otimes K)\preceq(Z\otimes K)^{T}(Z\otimes K)$,
whence $\aA\preceq P$. This is \eqref{eq:sandwich}.

Because $P\succ0$ (Lemma~\ref{lem:precond}(ii)) the symmetric matrix
$\widehat A:=P^{-1/2}\aA P^{-1/2}$ is well defined and similar to $P^{-1}\aA$,
so they share the spectrum $\{\mu\}$, which equals the generalized spectrum of
$(\aA,P)$. Pre- and post-multiplying \eqref{eq:sandwich} by $P^{-1/2}$ gives
$P^{-1/2}\lambda(I_r\otimes K)P^{-1/2}\preceq\widehat A\preceq I_{nr}$, so by
the Rayleigh--Ritz characterization every eigenvalue $\mu$ of $\widehat A$
satisfies $\mu\le1$ and
$\mu\ge\lambda_{\min}\!\big(P^{-1/2}(I_r\otimes K)P^{-1/2}\big)\,\lambda=\kappa_0$.
Since $\lambda>0$ and $I_r\otimes K\succ0$, the matrix
$P^{-1/2}\lambda(I_r\otimes K)P^{-1/2}\succ0$, so $\kappa_0>0$; this is
\eqref{eq:specbound}. The condition-number bound is immediate, and $\kappa_0$
depends only on $K$, $G$ and $\lambda$ (through $P$), not on the number of
observations $q$ nor on the sampling pattern $S$.
\end{proof}

\begin{remark}[Why this preconditioner]\label{rem:whyprecond}
$P$ in \eqref{eq:Pdef} is the operator one would solve if every entry were
observed; it captures the regularizer $\lambda(I_r\otimes K)$ exactly and the
data term up to the sampling projector $SS^{T}$. By Lemma~\ref{lem:specequiv}
the preconditioned spectrum lies in $[\kappa_0,1]$, so the standard CG bound
gives, after $m$ steps, a residual reduction
$\|w-w_m\|_{\aA}\le2\big(\tfrac{1-\sqrt{\kappa_0}}{1+\sqrt{\kappa_0}}\big)^{m}\|w-w_0\|_{\aA}$,
i.e.\ an iteration count $m=O(\kappa_0^{-1/2}\log(1/\epsilon))$ that is
\emph{independent of $q$}. This is the precise sense in which $P$ clusters the
spectrum and controls the PCG iteration count.
\end{remark}

\begin{lemma}[Right-hand side without order $M$ or $N$]\label{lem:rhs}
The right-hand side $b=(I_r\otimes K)\vecop(B)=\vecop(KB)$ with $B=TZ$ is
computed as follows, avoiding all order-$M$ and order-$N$ work:
\begin{enumerate}[label=\textup{(\roman*)}]
\item the sparse MTTKRP $B=TZ\in\bR^{n\times r}$ is
      $B=\sum_{j=1}^{q}t_j\,e_{a_j}\,\tilde Z_{j,:}$, i.e.\
      $B_{a,:}=\sum_{j:\,a_j=a} t_j\,\tilde Z_{j,:}$ where $t_j$ is the value of
      the $j$-th observed entry, at cost $O(qr)$ (or $O(qdr)$ including the
      row-of-$Z$ assembly of \eqref{eq:tilde}); if $B$ is supplied as input it
      is free;
\item $KB$ ($n\times n$ times $n\times r$) costs $O(n^{2}r)$.
\end{enumerate}
Hence forming $b$ costs $O(qr+n^{2}r)$, with no $O(M)$ or $O(N)$ term.
\end{lemma}

\begin{proof}
(i) $T$ has nonzeros only at the observed positions $(a_j,c_j)$ with values
$t_j$, so $T=\sum_j t_j e_{a_j}e_{c_j}^{T}$ and
$B=TZ=\sum_j t_j e_{a_j}(e_{c_j}^{T}Z)=\sum_j t_j e_{a_j}\tilde Z_{j,:}$,
which is the stated sparse accumulation: each $j$ adds the $1\times r$ vector
$t_j\tilde Z_{j,:}$ to row $a_j$ of $B$, costing $O(r)$ per observation, $O(qr)$
total (the row $\tilde Z_{j,:}$ itself was assembled in $O((d-1)r)$ in
\eqref{eq:tilde}, giving $O(qdr)$ if not cached). No dense $T\in\bR^{n\times M}$
or $Z\in\bR^{M\times r}$ is formed. (ii) The vec rule gives
$(I_r\otimes K)\vecop(B)=\vecop(KB)$, and $KB$ for dense $K$ costs $O(n^{2}r)$.
Adding (i) and (ii) gives $O(qr+n^{2}r)$.
\end{proof}

\section*{Algorithms}

Before stating the two required algorithm boxes we record the order-$M$/order-$N$
free assembly of $\aA$ used by the direct method.

\begin{lemma}[Observed-contraction assembly of $\aA$]\label{lem:assembly}
With the observed rows $\tilde K\in\bR^{q\times n}$, $\tilde Z\in\bR^{q\times r}$
of \eqref{eq:tilde}, the operator $\aA$ in \eqref{eq:Adef} equals
\begin{equation}\label{eq:contraction}
  \aA
  \;=\;\sum_{j=1}^{q}\bigl(\tilde Z_{j,:}^{T}\tilde Z_{j,:}\bigr)\otimes
        \bigl(\tilde K_{j,:}^{T}\tilde K_{j,:}\bigr)
      \;+\;\lambda\,(I_r\otimes K),
\end{equation}
where $\tilde Z_{j,:}^{T}\tilde Z_{j,:}\in\bR^{r\times r}$ and
$\tilde K_{j,:}^{T}\tilde K_{j,:}\in\bR^{n\times n}$ are the rank-one outer
products of the $j$-th observed rows. The sum is computed using only the
$q\times n$ and $q\times r$ arrays $\tilde K,\tilde Z$ and the $n\times n$ kernel
$K$; \emph{no $N$- or $M$-dimensional object is formed}, and the assembly costs
$O(qn^{2}r^{2})$.
\end{lemma}

\begin{proof}
Since $S$ is a column subset of $I_N$, $SS^{T}=\sum_{j=1}^{q}e_{p_j}e_{p_j}^{T}$
where $p_j=a_j+n\,c_j$ is the linear index of the $j$-th observation in
$\vecop(T)$ (so $e_{p_j}=e_{c_j}\otimes e_{a_j}$ under the column-major vec).
Hence
\[
  (Z\otimes K)^{T}SS^{T}(Z\otimes K)
  =\sum_{j=1}^{q}\bigl[(Z\otimes K)^{T}e_{p_j}\bigr]\bigl[(Z\otimes K)^{T}e_{p_j}\bigr]^{T}.
\]
Using $e_{p_j}=e_{c_j}\otimes e_{a_j}$ and the mixed-product rule,
$(Z\otimes K)^{T}(e_{c_j}\otimes e_{a_j})=(Z^{T}e_{c_j})\otimes(K^{T}e_{a_j})
=Z_{c_j,:}^{T}\otimes K_{a_j,:}^{T}=\tilde Z_{j,:}^{T}\otimes\tilde K_{j,:}^{T}$
(by symmetry $K_{a_j,:}^{T}=K_{:,a_j}$, and \eqref{eq:tilde}). Therefore the
$j$-th term is
$(\tilde Z_{j,:}^{T}\otimes\tilde K_{j,:}^{T})(\tilde Z_{j,:}^{T}\otimes\tilde K_{j,:}^{T})^{T}
=(\tilde Z_{j,:}^{T}\tilde Z_{j,:})\otimes(\tilde K_{j,:}^{T}\tilde K_{j,:})$,
again by the mixed-product rule. Adding $\lambda(I_r\otimes K)$ gives
\eqref{eq:contraction}.

\emph{Cost.} For each $j$: the outer product $\tilde Z_{j,:}^{T}\tilde Z_{j,:}$
is $O(r^{2})$, $\tilde K_{j,:}^{T}\tilde K_{j,:}$ is $O(n^{2})$, and the
Kronecker product of an $r\times r$ by an $n\times n$ matrix, accumulated into
$\aA$, is $O(n^{2}r^{2})$; summing over $j=1,\dots,q$ gives $O(qn^{2}r^{2})$.
Adding the block-diagonal $\lambda(I_r\otimes K)$ ($r$ copies of $K$) is
$O(n^{2}r)$, absorbed. Every quantity used has size $q$, $n$, or $r$; the index
ranges $M=\prod_{i\neq k}n_i$ and $N=\prod_i n_i$ never appear.
\end{proof}

\begin{remark}[Why the boxes are order-$M$/order-$N$ free]\label{rem:boxfree}
Both algorithm boxes below take as input \emph{only} the kernel $K$
($n\times n$), the per-observation triples $\{(a_j,c_j,t_j)\}_{j=1}^{q}$ (size
$O(q)$), and the fixed factors $\{A_i\}_{i\neq k}$ (sizes $n_i\times r$, total
$O(r\sum_{i\neq k}n_i)$, never their Khatri--Rao product $Z\in\bR^{M\times r}$).
The dense unfolding $T\in\bR^{n\times M}$ and the vectorization
$\vecop(T)\in\bR^{N}$ are \emph{never} materialized; $T$ enters only through the
$q$ stored values $t_j$. The Khatri--Rao matrix $Z\in\bR^{M\times r}$ is
\emph{never} materialized; it enters only through its $q$ observed rows
$\tilde Z$, each assembled by \eqref{eq:tilde} as a Hadamard product of $d-1$
factor rows in $O(dr)$. Consequently every line of both boxes carries a cost
bound in terms of $q,n,r,d$ only, with no $M$ or $N$ factor, as annotated.
\end{remark}

\begin{algorithm}[h]
\caption{Direct symmetric (Cholesky) solve of the mode-$k$ subproblem}
\label{alg:direct}
\begin{algorithmic}[1]
\Require Kernel $K\succ0$ ($n\times n$); per-observation triples
  $\{(a_j,c_j,t_j)\}_{j=1}^q$; fixed factors $\{A_i\}_{i\neq k}$; $\lambda>0$
  \Comment{inputs of size $O(n^2+q+r\sum_{i\neq k}n_i)$; no $Z\in\bR^{M\times r}$, no $T\in\bR^{n\times M}$}
\Ensure $W\in\bR^{n\times r}$ solving $\aA\,\vecop(W)=b$
\State $\tilde K\gets[K_{a_j,:}]_{j=1}^q\in\bR^{q\times n}$ \Comment{read $q$ rows of $K$: $O(qn)$}
\State $\tilde Z\gets[\,\circledast_{i\neq k}(A_i)_{c_j^{(i)},:}\,]_{j=1}^q\in\bR^{q\times r}$
  \Comment{Eq.~\eqref{eq:tilde}, Hadamard of $d-1$ factor rows: $O(qdr)$}
\State $\aA\gets\sum_{j=1}^q(\tilde Z_{j,:}^{T}\tilde Z_{j,:})\otimes(\tilde K_{j,:}^{T}\tilde K_{j,:})+\lambda(I_r\otimes K)$
  \Comment{Lemma~\ref{lem:assembly}: $O(qn^2r^2)$, no $O(M),O(N)$}
\State $B\gets\sum_{j=1}^q t_j\,e_{a_j}\tilde Z_{j,:}$ (sparse MTTKRP);\ \ $b\gets\vecop(KB)$
  \Comment{Lemma~\ref{lem:rhs}: $O(qr+n^2r)$, no $O(M),O(N)$}
\State Cholesky factor $\aA=R^{T}R$ ($R$ upper triangular, $nr\times nr$)
  \Comment{$O((nr)^3)=O(n^3r^3)$}
\State Solve $R^{T}u=b$ then $R\,\vecop(W)=u$ by triangular substitution
  \Comment{$O(n^2r^2)$}
\State \Return $W$
\end{algorithmic}
\end{algorithm}

\noindent
\emph{Direct-box accounting.} Lines 1--4 are the order-$M$/$N$-free assembly:
$O(qn+qdr)$ for the observed rows, $O(qn^2r^2)$ for $\aA$
(Lemma~\ref{lem:assembly}), $O(qr+n^2r)$ for $b$ (Lemma~\ref{lem:rhs}); lines
5--6 are the symmetric solve $O(n^3r^3)+O(n^2r^2)$. Total
$O(qn^2r^2+n^3r^3)$; we keep both terms rather than assuming one dominates, and
under $q<N$, $n\le M$ neither term involves $M$ or $N$. The dominant
term is the Cholesky factorization $O(n^3r^3)$, matching the problem statement's
cost for a standard symmetric solver, but now \emph{without} the forbidden
$O(M)$/$O(N)$ explicit-formation cost.

\begin{algorithm}[h]
\caption{Preconditioned CG (PCG) for the mode-$k$ subproblem}
\label{alg:pcg}
\begin{algorithmic}[1]
\Require Kernel $K\succ0$ ($n\times n$); per-observation triples
  $\{(a_j,c_j,t_j)\}_{j=1}^q$; fixed factors $\{A_i\}_{i\neq k}$; $\lambda>0$;
  tolerance $\epsilon\in(0,1)$; iteration cap $m_{\max}$
  \Comment{no $Z\in\bR^{M\times r}$, no $T\in\bR^{n\times M}$}
\Ensure $W\in\bR^{n\times r}$ with $\|\aA\,\vecop(W)-b\|_2\le\epsilon\|b\|_2$ (or $m_{\max}$ reached)
\Statex \textbf{One-time setup} (all lines $O(M),O(N)$-free)
\State $\tilde K\gets[K_{a_j,:}]_{j=1}^q$;\ \ $\tilde Z\gets[\circledast_{i\neq k}(A_i)_{c_j^{(i)},:}]_{j=1}^q$
  \Comment{Eq.~\eqref{eq:tilde}: $O(qn+qdr)$}
\State $G\gets\circledast_{i\neq k}(A_i^{T}A_i)\in\bR^{r\times r}$
  \Comment{Lemma~\ref{lem:precond}(i): $O(r^2\sum_{i\neq k}n_i+dr^2)$, no $O(M)$}
\State $K=U\Theta U^{T}$ ($O(n^3)$);\ \ $G=V\Gamma V^{T}$ ($O(r^3)$);\ \
  $\Omega_{a,p}\gets g_p\theta_a^2+\lambda\theta_a$ \Comment{one-time eig.; $O(nr)$ for $\Omega$}
\State $B\gets\sum_{j=1}^q t_j\,e_{a_j}\tilde Z_{j,:}$;\ \ $b\gets\vecop(KB)$
  \Comment{Lemma~\ref{lem:rhs}: $O(qr+n^2r)$, no $O(M),O(N)$}
\Statex \textbf{PCG iteration} (all lines $O(M),O(N)$-free): $\aA$-action by Lemma~\ref{lem:matvec} (Eq.~\eqref{eq:Aaction}); $P^{-1}$ by Eq.~\eqref{eq:Pinv}
\State $w\gets 0_{n\times r}$;\ \ $\rho\gets b$ (residual, since $\aA w=0$);\ \
  $z\gets P^{-1}\rho$ via \eqref{eq:Pinv};\ \ $p\gets z$;\ \ $\gamma\gets\rho^{T}z$
  \Comment{$O(n^2r+nr^2)$}
\For{$\ell=0,1,\dots,m_{\max}-1$}
  \State $v\gets\aA\,p$ via \eqref{eq:Aaction}
    \Comment{Lemma~\ref{lem:matvec}: $O(qnr+n^2r)$}
  \State $\alpha\gets\gamma/(p^{T}v)$ \Comment{$O(nr)$}
  \State $w\gets w+\alpha p$;\ \ $\rho\gets\rho-\alpha v$ \Comment{$O(nr)$}
  \If{$\|\rho\|_2\le\epsilon\|b\|_2$} \textbf{break} \EndIf \Comment{stopping test, $O(nr)$}
  \State $z\gets P^{-1}\rho$ via \eqref{eq:Pinv} \Comment{$O(n^2r+nr^2)$}
  \State $\gamma^{+}\gets\rho^{T}z$;\ \ $\beta\gets\gamma^{+}/\gamma$;\ \ $\gamma\gets\gamma^{+}$
    \Comment{$O(nr)$}
  \State $p\gets z+\beta p$ \Comment{$O(nr)$}
\EndFor
\State \Return $W\gets w$
\end{algorithmic}
\end{algorithm}

\noindent
\emph{PCG-box accounting and guarantees.} Every setup line (1--4) and every
iteration line (7--14) carries a cost bound in $q,n,r,d$ only; no line forms
$Z\otimes K$ ($nr\times nr$), the dense $T$ ($n\times M$), $\vecop(T)$ ($N$), or
$Z$ ($M\times r$). The vectors $w,\rho,z,p$ are stored as $n\times r$ matrices
($O(nr)$ memory). The \emph{stopping criterion} is the standard relative
residual $\|\rho_\ell\|_2\le\epsilon\|b\|_2$ (line 10); because $w_0=0$ this is
$\|\aA w_\ell-b\|_2\le\epsilon\|b\|_2$. The \emph{iteration bound} is two-fold:
(a) in exact arithmetic CG terminates in at most $\operatorname{rank}(\aA)\le nr$
steps (Lemma~\ref{lem:complexity}); (b) the preconditioned condition number
satisfies $\kappa_P(\aA)\le1/\kappa_0$ (Lemma~\ref{lem:specequiv}), so the
standard CG error bound gives
$\|w-w_\ell\|_{\aA}\le2\bigl(\tfrac{1-\sqrt{\kappa_0}}{1+\sqrt{\kappa_0}}\bigr)^{\ell}\|w-w_0\|_{\aA}$
and the tolerance $\epsilon$ is reached after
$m=\bigl\lceil\tfrac12\kappa_0^{-1/2}\log(2/\epsilon)\bigr\rceil$ iterations,
\emph{independent of $q$} (Remark~\ref{rem:whyprecond}). Hence one takes
$m_{\max}=\min\{nr,\lceil\tfrac12\kappa_0^{-1/2}\log(2/\epsilon)\rceil\}$ as a
safe a priori cap. The exact running time is tallied in
Lemma~\ref{lem:complexity}.

\section*{Complexity (Lemma L5)}

\begin{lemma}[Exact PCG complexity and comparison to the direct solve]
\label{lem:complexity}
Let $m$ be the number of PCG iterations performed by
Algorithm~\ref{alg:pcg}. Then:
\begin{enumerate}[label=\textup{(\roman*)}]
\item \textbf{Per iteration.} One PCG step costs one $\aA$-multiply
      $O(qnr+n^{2}r)$ (Lemma~\ref{lem:matvec}), one preconditioner solve
      $O(n^{2}r+nr^{2})$ (Lemma~\ref{lem:precond}(ii)), and $O(nr)$ for the CG
      vector updates (axpys/inner products on $nr$-vectors), so
      \begin{equation}\label{eq:periter}
        C_{\mathrm{iter}}=O\!\big(qnr+n^{2}r+nr^{2}\big)=O(qnr)\quad(\text{since }q>n,r).
      \end{equation}
\item \textbf{One-time setup.} The setup costs of Algorithm~\ref{alg:pcg} are
      \begin{equation}\label{eq:onetime}
        C_{\mathrm{setup}}
        = O\!\big(\underbrace{n^{3}}_{\text{eig }K}
        +\underbrace{r^{3}}_{\text{eig }G}
        +\underbrace{r^{2}\textstyle\sum_{i\neq k}n_i+dr^{2}}_{\text{Gram }G}
        +\underbrace{qn+qdr}_{\tilde K,\tilde Z}
        +\underbrace{qr+n^{2}r}_{\text{RHS}}\big)
        = O\!\big(n^{3}+r^{3}+qdr\big),
      \end{equation}
      where the last simplification uses $n,r<q$ and
      $\sum_{i\neq k}n_i\le (d-1)\max_i n_i< (d-1)q$.
\item \textbf{Total.}
      \begin{equation}\label{eq:total}
        C_{\mathrm{PCG}}
        = m\cdot C_{\mathrm{iter}}+C_{\mathrm{setup}}
        = O\!\big(m\,qnr\;+\;n^{3}+r^{3}+qdr\big).
      \end{equation}
\item \textbf{Comparison to the direct $O(n^{3}r^{3})$ solve.} The ratio is
      \begin{equation}\label{eq:ratio}
        \frac{C_{\mathrm{PCG}}}{n^{3}r^{3}}
        = O\!\Big(\frac{m\,q}{n^{2}r^{2}}+\frac1{r^{3}}+\frac1{n^{3}}
        +\frac{q d}{n^{3}r^{2}}\Big),
      \end{equation}
      so PCG is asymptotically cheaper than the direct solve whenever
      $m\,q=o(n^{2}r^{2})$ (the one-time terms are already
      $o(n^{3}r^{3})$ because $n,r\ge2$, $d<n$ and $q<n^{2}r^{2}$ in any regime
      with $q<N$ and $n^2\le N$). In the absolute worst case CG terminates in
      $m\le\operatorname{rank}(\aA)\le nr$ exact-arithmetic steps, giving
      $C_{\mathrm{PCG}}=O(qn^{2}r^{2}+n^{3}+r^{3})$, smaller than $n^{3}r^{3}$
      by the factor $\Theta\!\big(q/(nr)\big)$ on the dominant term; this is an
      improvement precisely when $q=o(nr)$. In practice $m\ll nr$ because the
      preconditioner clusters the spectrum (Remark~\ref{rem:whyprecond}), so
      the realized cost $O(m\,qnr)$ is far below both $qn^2r^2$ and $n^3r^3$.
\end{enumerate}
\end{lemma}

\begin{proof}
(i) Sum the per-step costs: the mat--vec $O(qnr+n^{2}r)$ from
\eqref{eq:matveccost}, the preconditioner apply $O(n^{2}r+nr^{2})$ from
Lemma~\ref{lem:precond}(ii), and the standard CG bookkeeping (two inner
products, three axpys on vectors of length $nr$) $O(nr)$. The maximum of these
is $O(qnr+n^{2}r+nr^{2})$; since $q>n$ gives $qnr\ge n^{2}r$ and $q>r$ gives
$qnr\ge nr^{2}$, the dominant term is $qnr$, proving \eqref{eq:periter}.

(ii) Add the setup lines of Algorithm~\ref{alg:pcg}: assembling $\tilde K$,
$\tilde Z$ is $O(qn+qdr)$ (read $q$ rows of $K$; build $q$ rows of $Z$ as
Hadamard products of $d-1$ factor rows); forming $G$ via \eqref{eq:krgram} is
$O(r^{2}\sum_{i\neq k}n_i+dr^{2})$; the eigendecompositions of $K$ and $G$ are
$O(n^{3})$ and $O(r^{3})$; the RHS is $O(qr+n^{2}r)$ by Lemma~\ref{lem:rhs}.
Using $n^{2}r\le qnr\le n^3+ \text{(absorbed)}$, more simply $n^2r<n^3$ as $r<q$
and $n^2 r\le n\cdot nr$, and $\sum_{i\neq k}n_i<(d-1)q$, all subdominant terms
are absorbed into $O(n^{3}+r^{3}+qdr)$, giving \eqref{eq:onetime}.

(iii) Immediate from $C_{\mathrm{PCG}}=m\,C_{\mathrm{iter}}+C_{\mathrm{setup}}$
and (i)--(ii).

(iv) Divide \eqref{eq:total} term-by-term by $n^{3}r^{3}$:
$m\,qnr/(n^{3}r^{3})=mq/(n^{2}r^{2})$, $n^{3}/(n^{3}r^{3})=1/r^{3}$,
$r^{3}/(n^{3}r^{3})=1/n^{3}$, and $qdr/(n^{3}r^{3})=qd/(n^{3}r^{2})$, which is
\eqref{eq:ratio}. The one-time terms are subdominant: $1/r^{3}\to0$ and
$1/n^{3}\to0$ as $n,r\to\infty$, while $d<n$ gives
$qd/(n^{3}r^{2})<q/(n^{2}r^{2})$, so this term is also $o(1)$ exactly when
$q=o(n^{2}r^{2})$. Thus, under the explicitly stated sufficient condition
$m\,q=o(n^{2}r^{2})$ (which subsumes $q=o(n^{2}r^{2})$ since $m\ge1$), every
term of \eqref{eq:ratio} tends to $0$ and PCG is asymptotically cheaper than the
direct solve. The worst-case bound uses the finite
termination of CG: in exact arithmetic CG/PCG converges in at most
$\operatorname{rank}(\aA)\le nr$ steps (the dimension of the Krylov space is at
most the number of distinct eigenvalues, itself $\le nr$), so $m\le nr$ and
$m\,qnr\le qn^{2}r^{2}$; dividing by $n^{3}r^{3}$ gives the factor $q/(nr)$ and
the stated improvement region $q=o(nr)$.
\end{proof}

\begin{remark}[Summary of the accounting]\label{rem:summary}
The PCG solver never forms $\aA$ (avoiding the $O(qn^{2}r^{2})$ assembly and
$O(n^{3}r^{3})$ factorization of the direct method), never touches $\vecop(T)$
or $Z$ in full (avoiding $O(N)$ and $O(M)$), uses the exact dense $K$-multiply
cost $O(n^{2}r)$ inside the $O(qnr)$ mat--vec, pays the one-time
$O(n^{3})$ Cholesky/eigendecomposition of $K$ and $O(r^{3})$ factorization of
the $r\times r$ Gram $G$ (both $\ll n^{3}r^{3}$), and forms the RHS
$\vecop(KB)$ from the precomputed/sparse MTTKRP $B=TZ$ at $O(qr+n^{2}r)$. The
total $O(m\,qnr+n^{3}+r^{3}+qdr)$ beats the direct $O(n^{3}r^{3})$ on the
dominant term by the factor $mq/(n^{2}r^{2})$, which is $\le q/(nr)$ even in the
worst case $m=nr$.
\end{remark}

\end{document}
